% Fichier complété par tools/complete_missing_corrections.py.
% Source : CIN/CIN-03-Transmetteurs/36_VisEcrou/36_VisEcrou.tex
% Statut : rédigé (3 question(s)).
\begin{corrigebox}[Corrigé rédigé]
\CorrigeQuestion{1}
Chaîne : moteur 1 $(C_m,\omega_1)$ $\rightarrow$ transmission par
                courroie/poulies $(\omega_2)$ $\rightarrow$ vis--écrou de pas $p_v$
                $\rightarrow$ piston 3 $(F_3,v_3)$.
\CorrigeQuestion{2}
Sans glissement, $\omega_2/\omega_1=D_1/D_2=1/2$. Pour la vis--écrou,
                $v_3=p_v\omega_2/(2\pi)$. Ainsi
                \[\boxed{v_3=\frac{p_v}{4\pi}\,\omega_1}.\]
\CorrigeQuestion{3}
Le sens de translation est donné par le pas de la vis et le sens de
                rotation. Le signe de la relation précédente est donc à adapter à la
                convention positive choisie; sa valeur absolue reste $p_v/(4\pi)$.
\end{corrigebox}
